% PAGES: 185-186
% Continues directly after sec06_c1_4.tex (pages 183-184). That file ended
% mid-display-math, inside an UNCLOSED "\[" whose last line was
% "\end{pmatrix}," -- the sentence "Then, (U(2:4,2:4))^t U(2:4,2:4) =
% [numeric matrix]," (folio 168) continues here with the closing "\]" below,
% followed immediately by "which can be written as" (folio 169). Do not
% reopen a "\[" before that "\]".
% This file also closes the "\begin{examplex}" opened in sec06_c1_4.tex
% (folio 168, "Example: Let $A$ be the following symmetric matrix:") via
% "\end{examplex}" below -- no "\begin{examplex}" or "Example:" label is
% reprinted here, and the environment's auto-appended closing square is the
% only square used (no manual "\square").
\]
which can be written as
\[
\begin{pmatrix} U(2,2) & 0 & 0 \\ U(2,3) & U(3,3) & 0 \\ U(2,4) & U(3,4) & U(4,4) \end{pmatrix}
\begin{pmatrix} U(2,2) & U(2,3) & U(2,4) \\ 0 & U(3,3) & U(3,4) \\ 0 & 0 & U(4,4) \end{pmatrix}
\]
\[
\;=\; \begin{pmatrix} 4 & -2 & 6 \\ -2 & 17 & 5 \\ 6 & 5 & 14 \end{pmatrix}.
\]
The unknown entries from the second row of $U$ can be computed as the first row of
the Cholesky factor of the $3 \times 3$ matrix above, as follows:
\[
U(2,2) \;=\; \sqrt{4} \;=\; 2;
\]
\[
U(2,3) \;=\; \frac{-2}{U(2,2)} \;=\; \frac{-2}{2} \;=\; -1; \quad
U(2,4) \;=\; \frac{6}{U(2,2)} \;=\; \frac{6}{2} \;=\; 3.
\]
The current form of $U$ is
\begin{equation}
U \;=\; \begin{pmatrix} 3 & -1 & 2 & -1 \\ 0 & 2 & -1 & 3 \\ 0 & 0 & U(3,3) & U(3,4) \\
0 & 0 & 0 & U(4,4) \end{pmatrix}.
\end{equation}
The updated form of the $2 \times 2$ matrix $A(3:4,3:4)$ is computed using (6.29), from
$A(3:4,3:4)$ obtained from (6.32) and with $U$ given by (6.33), as follows:
\begin{align*}
A(3:4,3:4)
&= A(3:4,3:4) \;-\; (U(2,3:4))\tp U(2,3:4) \\
&= \begin{pmatrix} 17 & 5 \\ 5 & 14 \end{pmatrix} \;-\;
\begin{pmatrix} -1 \\ 3 \end{pmatrix} \begin{pmatrix} -1 & 3 \end{pmatrix} \\
&= \begin{pmatrix} 17 & 5 \\ 5 & 14 \end{pmatrix} \;-\;
\begin{pmatrix} 1 & -3 \\ -3 & 9 \end{pmatrix} \\
&= \begin{pmatrix} 16 & 8 \\ 8 & 5 \end{pmatrix}.
\end{align*}
Thus, the updated form of the $2 \times 2$ matrix $A(3:4,3:4)$ is
\begin{equation}
A(3:4,3:4) \;=\; \begin{pmatrix} 16 & 8 \\ 8 & 5 \end{pmatrix}.
\end{equation}
Then,
\[
(U(3:4,3:4))\tp U(3:4,3:4) \;=\; \begin{pmatrix} 16 & 8 \\ 8 & 5 \end{pmatrix},
\]
which can be written as
\[
\begin{pmatrix} U(3,3) & 0 \\ U(3,4) & U(4,4) \end{pmatrix}
\begin{pmatrix} U(3,3) & U(3,4) \\ 0 & U(4,4) \end{pmatrix}
\;=\; \begin{pmatrix} 16 & 8 \\ 8 & 5 \end{pmatrix}.
\]
The unknown entries from the third row of $U$ can be computed as the first row of
the Cholesky factor of the $2 \times 2$ matrix above, as follows:
\[
U(3,3) \;=\; \sqrt{16} \;=\; 4; \quad U(3,4) \;=\; \frac{8}{U(3,3)} \;=\; \frac{8}{4} \;=\; 2.
\]

The current form of $U$ is
\begin{equation}
U \;=\; \begin{pmatrix} 3 & -1 & 2 & -1 \\ 0 & 2 & -1 & 3 \\ 0 & 0 & 4 & 2 \\
0 & 0 & 0 & U(4,4) \end{pmatrix}.
\end{equation}
The updated form of $A(4,4)$, which is a number, is computed using (6.29), from
$A(4,4)$ obtained from (6.34) and with $U$ given by (6.35), as follows:
\[
A(4,4) \;=\; A(4,4) - U(3,4)U(3,4) \;=\; 5 - 2 \cdot 2 \;=\; 1,
\]
which corresponds to
\[
(U(4,4))^2 \;=\; 1.
\]
Thus, $U(4,4) = \sqrt{1} = 1$ since $U(4,4)$ must be positive. We conclude that $A$ is a
symmetric positive definite matrix with the following Cholesky factor:
\[
U \;=\; \begin{pmatrix} 3 & -1 & 2 & -1 \\ 0 & 2 & -1 & 3 \\ 0 & 0 & 4 & 2 \\
0 & 0 & 0 & 1 \end{pmatrix}.
\]
\end{examplex}

The pseudocode for the Cholesky decomposition can be found in Table 6.1.

\begin{table}[H]
\centering
\caption{Pseudocode for Cholesky decomposition}
\fbox{\begin{minipage}{0.85\linewidth}
\textbf{Function Call:}\\
$U = \text{cholesky}(A)$
\bigskip

\textbf{Input:}\\
$A = $ symmetric positive definite matrix of size $n$
\bigskip

\textbf{Output:}\\
$U = $ upper triangular matrix such that $U\tp U = A$
\bigskip

\noindent
for $i = 1:(n-1)$\\
\hspace*{1.5em}$U(i,i) = \sqrt{A(i,i)}$;\\
\hspace*{1.5em}for $k = (i+1):n$\\
\hspace*{3em}$U(i,k) = A(i,k)/U(i,i)$; \hfill\textit{// compute row $i$ of $U$}\\
\hspace*{1.5em}end\\
\hspace*{1.5em}for $j = (i+1):n$\\
\hspace*{3em}for $k = j:n$\\
\hspace*{4.5em}$A(j,k) = A(j,k) - U(i,j)U(i,k)$;\\
\hspace*{3em}end\\
\hspace*{1.5em}end\\
end\\
$U(n,n) = \sqrt{A(n,n)}$
\end{minipage}}
\end{table}

Recall from Lemma 2.3 that the operation count for the LU decomposition is
$\frac{2}{3}n^3 + O(n^2)$. As expected, the operation count for the Cholesky
decomposition is approximately half of the operation count for the LU decomposition:
